洗澡时先把冷热水龙头拧到昨天试出的刻度，然后闭眼站进去——这是开环：动作按预先定好的表执行。伸手试水温，凉了拧热、烫了拧冷——这是闭环：动作由当前观测决定。两者的差别与算得精不精细无关，只看新的动作有没有用上刚刚看到的结果。模型完全准确、世界完全安静时，两者给出同一条轨迹；而这两个前提从不成立，反馈的全部价值就压在这句话上。

本单元先站在最强的假设上：动力学 \(\dot x = f(x,u)\) 与代价函数都已知，问题只剩把``控制得好''写成数学并解出来；最后两篇再撤掉``完整状态可见''与``模型天然已知''这两项。六篇的顺序是从反馈直觉走到最优性条件，再走到可计算的算法，最后回到数据。

\begin{enumerate}
\tightlist
\item
  从开环到闭环反馈：为什么事先算好的计划对模型误差与扰动脆弱，反馈怎样把闭环系统改造成以目标为汇的梯度流，PID 三项各治什么病、各带什么副作用。
\item
  最大值原理：扰动一条候选最优轨线做一阶变分，引入协态消掉状态偏移，把无穷维寻优压成每一时刻的逐点选择，代价是问题落成一个两点边值问题。
\item
  HJB 方程与动态规划：把``从任意状态出发的最优剩余代价''定义成值函数，得到 Hamilton--Jacobi--Bellman 方程，在 LQR 上完整算到 Riccati 方程，并撞上维数灾难。
\item
  LQG：线性动力学、二次代价、Gaussian 噪声三条同时成立时，最优输出反馈可以拆成 Kalman 滤波与 LQR 增益两块独立设计——以及这个分离在什么地方不再成立。
\item
  DDP：沿一条名义轨迹展开动力学与代价，反向递推得到局部反馈律，前向 rollout 用真实动力学检验，它是 Newton 法搬到轨迹空间的版本。
\item
  系统辨识：动力学从输入输出数据估出来，于是继承统计推断的全部问题——持续激励、闭环相关、状态坐标的不唯一，以及辨识误差在闭环里被放大的机制。
\end{enumerate}

两条主线值得先记住。最大值原理与 HJB 是同一个问题的两个视角，前者给必要条件，后者给动态规划的连续时间形式，沿最优轨线上协态恰好等于值函数的梯度；两者的账单形态不同，一个是 \(2n\) 维两点边值问题，一个是 \(n+1\) 维偏微分方程。另一条线是假设的逐级松动：LQR 靠线性加二次换来 Riccati 递推，LQG 靠高斯换来分离，DDP 放弃全局只求局部改进，系统辨识连模型本身都交给数据。松一格，保证就弱一格。

记号在全单元统一：状态记 \(x\)，控制输入记 \(u\)，值函数记 \(V\)，协态记 \(\lambda\)。

\subsection{怎么读这一章}

首次阅读抓住三层：反馈、目标、价值。第一篇建立反馈直觉，回答``为什么必须看着现在做决定''；第二与第三篇是同一枚硬币的两面，一面是一条候选轨线上的必要条件，一面是所有状态上的最优剩余代价，读完应当能说清两者各自适合什么场景。后三篇是三种结构化的出路，分别依赖线性加高斯、局部二次近似、有信息量的数据。

公式可以后置，假设不能后置。每一篇的推导都挂在几条前提上，前提被现实拿走一条，结论就换一副面孔——把这些前提读进去，比把公式抄下来有用得多。

\subsection{本章篇目}

以下篇目按概念依赖排列；若从具体问题进入，也可以先打开最接近当前困惑的一篇，再沿篇内的链接回补。

\begin{itemize}
\tightlist
\item \hyperref[sec:16-Control-and-Reinforcement-Learning/01-Optimal-Control/01]{``从开环到闭环反馈''}；
\item \hyperref[sec:16-Control-and-Reinforcement-Learning/01-Optimal-Control/02]{``最大值原理：把最优轨线变成两点边值''}；
\item \hyperref[sec:16-Control-and-Reinforcement-Learning/01-Optimal-Control/03]{``HJB 方程与动态规划''}；
\item \hyperref[sec:16-Control-and-Reinforcement-Learning/01-Optimal-Control/04]{``Kalman 滤波：用预测与更新的循环从噪声中估计状态''}；
\item \hyperref[sec:16-Control-and-Reinforcement-Learning/01-Optimal-Control/05]{``LQG：看不全状态时怎样继续反馈''}；
\item \hyperref[sec:16-Control-and-Reinforcement-Learning/01-Optimal-Control/06]{``DDP：用局部二次模型改进整条轨迹''}；
\item \hyperref[sec:16-Control-and-Reinforcement-Learning/01-Optimal-Control/07]{``系统辨识：从输入输出数据估计动力学''}。
\end{itemize}

读完这七篇，剩下的缺口是模型本身建不起来的情形：状态空间大到画不出地图，动力学复杂到写不成方程，只能靠与环境反复交互攒经验。\hyperref[ch:136]{``强化学习：模型未知时怎样逼近最优决策''} 从那里接手，而它要绕开的东西，正是本单元第三篇撞上的维数灾难。
